%\documentclass[a4paper]{article}
\documentclass[12pt]{article}
%\documentclass[11pt]{article}

\usepackage[utf8]{inputenc}   % LaTeX, comprends les accents !
\usepackage[T1]{fontenc}      % Police contenant les caract�res fran�ais
\usepackage[francais]{babel}  % Placez ici une liste de langues, la
                              % derni�re �tant la langue principale

\usepackage{amsfonts}         % Augmente le nombre de symboles math�matiques accessibles.
\usepackage{amssymb}
\usepackage{amsmath}

\usepackage[a4paper]{geometry}% R�duire les marges
% \pagestyle{headings}        % Pour mettre des ent�tes avec les titres
                              % des sections en haut de page

\author{Woessner Guillaume}
\title{Révisions AL}         % Les param�tres du titre : titre, auteur, date
\date{}                       % La date n'est pas requise (la date du
                              % jour de compilation est utilis�e en son
			      % absence)

\sloppy                       % Ne pas faire d�border les lignes dans la marge

\begin{document}

\maketitle                % Faire un titre utilisant les donn�es

\paragraph{Contexte}~\\

\noindent On travaille sur les applications linéaires suivantes.

\[
\begin{array}{rccl}
f : & \mathbb{R}^3 & \longrightarrow & \mathbb{R}^3 \\
~& (x,y,z) & \longmapsto & (x+2y+3z ~;~ -2x-2z ~;~ -3x+2y-z)\\
~&~&~&\\
g : & \mathbb{R}_2[X] & \longrightarrow & \mathbb{R}^2 \\
~& P & \longmapsto & \Big((XP)(2) ~;~ (XP)'(2)\Big)\\
~&~&~&\\
h : & \mathbb{R}_2[X] & \longrightarrow & \mathcal{M}_{2\times 2}(\mathbb{R}) \\
~& P & \longmapsto & 
\begin{pmatrix}
P(0) & P'(0) \\
P'(0) & P(0) + P'(0)
\end{pmatrix}\\
~&~&~&\\
i : & \mathbb{R}_2[X] & \longrightarrow & \mathbb{R}_2[X] \\
~& P & \longmapsto & P - \frac{P''(1)}{2}
\end{array}
\]
~\\
Soient les bases alternatives suivantes.
\begin{itemize}
\item[$\bullet$] dans $\mathbb{R}^2$ : $\mathcal{B}_1 = \Big( (1,1) ~;~(1,-1)\Big)$.
\item[$\bullet$] dans $\mathbb{R}^3$ : $\mathcal{B}_2 = \Big( (1,1,1) ~;~ (1,-1,0) ~;~ (2,0,0)\Big)$.
\item[$\bullet$] dans $\mathbb{R}_2[X]$ : $\mathcal{B}_3 = \Big( X^2-1 ~;~ X-1 ~;~ -1\Big)$.
\end{itemize}
~\\
Soient les vecteurs-exemples suivants.
\begin{itemize}
\item[$\bullet$] dans $\mathbb{R}^2$ : $u=(4,-5)$.
\item[$\bullet$] dans $\mathbb{R}^3$ : $v=(3,0,-3)$.
\item[$\bullet$] dans $\mathbb{R}_2[X]$ : $w = 1 + 2X + X^2$.
%\item[$\bullet$] dans $\mathcal{M}_{2\times 2}(\mathbb{R})$ : $m = \begin{pmatrix}
%1 & -1 \\
%-2 & 2
%\end{pmatrix}$.
\end{itemize}

\newpage

\paragraph{Questions}~\\

\begin{itemize}
\item[1.] Trouver les coordonnées des vecteurs-exemples dans leur base altenative.\\
\item[2.] Quels sont les vecteurs de coordonnées $(1,2)$ (resp. $(1,2,3)$) dans les bases alternatives des espaces de dimension 2 (resp. de dimension 3) ?\\
\item[3.] Trouver le rang, la dimension du noyau de chacune de ces applications linéaires et trouver une base pour l'image et le noyau de chacune des applications linéaires.\\
\item[4.] Trouver la matrice de chacune des applications linéaires, dans la base canonique au départ et à l'arrivée.\\
\item[5.] Trouver la matrice de l'application $f$ dans la base alternative au départ, et la base canonique à l'arrivée.\\
\item[6.] Trouver la matrice de l'application $g$ dans la base canonique au départ, et la base alternative à l'arrivée.\\
\item[7.] Trouver la matrice de l'application $i$ dans la base alternative au départ, et la base alternative à l'arrivée.\\
\item[8.] Les applications sont-elles des projecteurs ?\\
\item[9.] Trouver le polynome caractéristique des endomorphismes.\\
\item[10.] Diagonaliser les endomorphismes.\\
\end{itemize}






\end{document}
